% Source: arithmetic/large_energy.tex




Put $R_m=\Z[z]/(z^m-1)$, let $z^*=z^{-1}$, and write
$J=1+z+\cdots+z^{m-1}$. We seek
\begin{equation}\label{s01:eq:target}
 A_i\equiv J\pmod{2R_m},\qquad
 \sum_{i=0}^3A_iA_i^*=4t\cdot1.
\end{equation}
When $m$ is odd, elementary energy and augmentation arguments force
$t\ge m$ and $t$ odd. At $t=m$, every coefficient is a sign, giving
four periodic complementary binary sequences and hence, by the
Goethals--Seidel construction, a Hadamard matrix of order $4m$.
The general boundary problem is not solved here.

\begin{theorem}[Uniform large-energy existence]\label{s01:thm:main}
If $m\ge1$ and $t$ are odd integers and $t\ge(m^2+3)/4$, then
\eqref{s01:eq:target} has a solution of the form
\[
 A_i=J+2x_i\cdot1,\qquad x_i\in\Z.
\]
For odd $m,t$, this particular form exists if and only if
$t\ge(m^2+3)/4$.
\end{theorem}
\begin{proof}
Since $JJ^*=mJ$ and $J^*=J$,
\begin{equation}\label{s01:eq:reduction}
 \sum_{i=0}^3(J+2x_i)(J+2x_i)^*
 =4\Bigl(m+\sum_i x_i\Bigr)J+4\sum_i x_i^2\cdot1.
\end{equation}
For $m>1$, $J$ and $1$ are linearly independent. Thus the target is
equivalent to
\begin{equation}\label{s01:eq:four}
 \sum_i x_i=-m,\qquad \sum_i x_i^2=t.
\end{equation}
For $m=1$, these conditions are still sufficient; necessity of the
stated threshold merely says $t\ge1$ and follows from coefficient energy.

Set $N=4t-m^2$. Under the hypotheses, $N\ge3$ and $N\equiv3\pmod8$.
The Legendre--Gauss three-square theorem gives
$N=u^2+v^2+w^2$ for integers $u,v,w$. All three are odd, because the
only square residues modulo eight are $0,1,4$. By changing the sign of
one of them if necessary, arrange
$u+v+w\equiv m\pmod4$. Define
\begin{align*}
 x_0&=(-m+u+v+w)/4,& x_1&=(-m+u-v-w)/4,\\
 x_2&=(-m-u+v-w)/4,& x_3&=(-m-u-v+w)/4.
\end{align*}
They are integers: the last three numerators differ from the first by
twice a sum of two odd numbers. Their sum is $-m$, and the orthogonal
four-square identity gives
\[
 \sum_i x_i^2=\frac{m^2+u^2+v^2+w^2}{4}=t.
\]
Equation~\eqref{s01:eq:reduction} proves existence.

Conversely, for $m>1$ any solution of this form obeys
\eqref{s01:eq:four}. Cauchy--Schwarz gives $4t\ge m^2$. Since
$4t-m^2\equiv3\pmod8$, it is at least three. This is precisely
$t\ge(m^2+3)/4$.
\end{proof}

The invoked classical three-square theorem has an elementary geometric
proof in Ankeny~\cite{s01:Ankeny}, whose Theorem 1 applies whenever the
integer is not $4^a(8b+7)$. No code or proof supplement from that paper
is used.

\begin{corollary}
The broader parity-constrained norm problem has a solution for every
odd $t\ge m$ when $m=1$ or $m=3$.
\end{corollary}
For $m=5$ the theorem covers every odd $t\ge7$; a single binary boundary
witness at $t=5$ therefore settles that entire fixed-$m$ question too.
It does not settle all $m$.

\begin{proposition}[Odd multiplier closure]\label{s01:prop:multiplier}
If \eqref{s01:eq:target} has a solution at $(m,t)$ and $s$ is any positive
odd integer, it has a solution at $(m,st)$.
\end{proposition}
\begin{proof}
By Lagrange's four-square theorem choose integers $a,b,c,d$ with
$a^2+b^2+c^2+d^2=s$. The integral matrix
\[
 M=\begin{pmatrix}
 a&b&c&d\\-b&a&-d&c\\-c&d&a&-b\\-d&-c&b&a
 \end{pmatrix}
\]
satisfies $M^{\mathsf T}M=sI_4$. Each row sum is odd, since its parity
equals $a+b+c+d\equiv s\equiv1\pmod2$.
Put $B_i=\sum_jM_{ij}A_j$. Then $B_i\equiv J\pmod2$, and expanding
the Hermitian norm sum using real integer coefficients of $M$ gives
\[
 \sum_i B_iB_i^*=s\sum_jA_jA_j^*=4st.
\]
This only uses a constant integral orthogonal design; it does not
assert a four-square composition identity with arbitrary group-ring
Hermitian coefficients.
\end{proof}

\section*{Examples and scope}
At $m=7$, the threshold is $t=13$. Taking $(u,v,w)=(1,1,1)$ gives
$(x_0,x_1,x_2,x_3)=(-1,-2,-2,-2)$, hence constant coefficients
$(-1,-3,-3,-3)$ and all remaining coefficients equal to one.
At $m=9$, the threshold is $21$; choosing $(u,v,w)=(-1,1,1)$ gives
$(-2,-3,-2,-2)$ and constant coefficients $(-3,-5,-3,-3)$.

For $m>3$, $(m^2+3)/4>m$, so the constant-perturbation family cannot
reach the Hadamard boundary. This follows from its sharp lower bound,
not from a failure to search far enough. Applying odd multiplier closure
to any known boundary solution gives infinitely many further energies,
but it cannot lower an energy to the boundary.

The remaining parameter strip for a fixed odd $m$ is finite:
\[
 m\le t<(m^2+3)/4,\qquad t\text{ odd}.
\]
It is an infinite family of strips as $m$ varies. A uniform theorem in
that strip, especially near $t=m$, is still needed. The polynomial
bound above is a proved partial answer to a stronger arithmetic question,
not a new Hadamard-order existence result.

\section*{A lower threshold from skew quadratic characters}
Suppose $m\equiv3\pmod4$ and a polynomial $L\in R_m$ has constant
coefficient zero, all other coefficients signs, and
\[
 L^*=-L,\qquad LL^*=m\cdot1-J.
\]
For every prime $p\equiv3\pmod4$ this is supplied by
$L=\sum_{j=1}^{p-1}\chi(j)z^j$, where $\chi$ is the quadratic character
modulo $p$, as in the Paley character construction~\cite{s01:Paley}.
For completeness, $\chi(-j)=-\chi(j)$ gives skewness. For $s\ne0$,
\[
 \sum_j\chi(j)\chi(j+s)
 =\sum_{j\ne0}\chi(1+s/j)=-1:
\]
as $j$ ranges over the nonzero field elements, $1+s/j$ ranges over all
field elements except one, and a nontrivial character sums to zero.
The constant coefficient is $p-1$, giving the norm identity.

\begin{theorem}[Skew-character large-energy region]\label{s01:thm:skew}
Under those hypotheses on $L$, the norm problem is solvable for every
odd
\[
 t\ge T_m:=\frac{m^2+6m+21}{16}.
\]
In particular this holds for every prime $m\equiv3\pmod4$.
The integer $T_m$ is odd. The threshold is sharp within the family
used in this proof.
\end{theorem}
\begin{proof}
Write $m=4k+3$ and $x=(3-m)/4=-k$. Then
$T_m=k^2+3k+3$ is an odd integer.
For odd integers $a,b,c$, set
\[
 A_1=L+a,\quad A_2=L+b,\quad A_3=L+c,\quad A_4=J+2x.
\]
Every coefficient is odd. Since $L^*=-L$, cross terms with a real
constant vanish, and
\[
 (L+a)(L+a)^*=(m+a^2)\cdot1-J.
\]
The fourth norm is $4x^2\cdot1+(m+4x)J=4x^2\cdot1+3J$.
Thus
\[
 \sum_{i=1}^4 A_iA_i^*
 =\bigl(3m+4x^2+a^2+b^2+c^2\bigr)\cdot1.
\]
The required integer for the last three squares is
$N=4t-3m-4x^2$. It is at least three precisely when $t\ge T_m$.
Also $3m+4x^2=4k^2+12k+9\equiv1\pmod8$, while $4t\equiv4\pmod8$;
hence $N\equiv3\pmod8$. The three-square theorem supplies odd
$a,b,c$ with squared sum $N$. Conversely, any three odd integers have
squared sum at least three, proving sharpness in this family.
\end{proof}

At $m=7$ this threshold is seven, so the entire broader problem is
settled for that fixed length as well. At $m=11$ it is thirteen, leaving
only the binary boundary case $t=11$ below the threshold. For large
primes the threshold remains quadratic in $m$, so it does not reach
the Hadamard boundary uniformly. The prime-power version of Paley's
character construction lives in the group ring of the additive field,
which need not be cyclic; no cyclic prime-power extension is asserted.

\begin{proposition}[A zero-aligned character ansatz is impossible]
With the same $m,L$, it is impossible that all four elements have the
form
\[
 A_i=u_iL+v_i+2w_iJ,\qquad u_i,v_i\text{ odd},\quad w_i\in\Z,
\]
and their Hermitian norm sum is a scalar multiple of $1$.
\end{proposition}
\begin{proof}
The coefficient multiplying $J$ in the sum of norms is
\[
 K=-\sum_i u_i^2+4\sum_i(w_iv_i+mw_i^2).
\]
Every summand in the last parentheses is even, since its parity is
$w_i+w_i^2=0$. The four odd squares sum to four modulo eight.
Therefore $K\equiv4\pmod8$, so it cannot vanish. The elements $J$ and
$1$ are linearly independent because $m>1$.
\end{proof}
This is an obstruction to a restricted character ansatz, not to the
general norm problem or to Hadamard existence. Theorem~\ref{s01:thm:skew}
avoids it by using an odd coefficient of $J$ in the fourth element.

\section*{Exact checks}
The standard-library script \path{verify_large_energy.py} constructs and
checks examples at the threshold and several larger odd energies for
every odd $m\le101$. It also tests odd multiplier closure using a direct
cyclic autocorrelation check. The general assertions rest on the paper
proofs; finite tests are additional checks only.

\begin{thebibliography}{2}
\bibitem{s01:Ankeny} N. C. Ankeny, \emph{Sums of three squares},
Proc. Amer. Math. Soc. 8 (1957), 316--319, Theorem 1.
\url{https://doi.org/10.1090/S0002-9939-1957-0085275-8}.
\bibitem{s01:Paley} R. E. A. C. Paley, \emph{On orthogonal matrices},
J. Math. Phys. 12 (1933), 311--320.
\url{https://doi.org/10.1002/sapm1933121311}.
\end{thebibliography}

